\documentclass[12pt,a4paper]{article}

\usepackage[utf8]{inputenc}
\usepackage[T1]{fontenc}
\usepackage{mathptmx}
\usepackage[a4paper,margin=1in]{geometry}
\usepackage{setspace}
\usepackage{titlesec}
\usepackage{fancyhdr}
\usepackage[hidelinks]{hyperref}
\usepackage{xurl}

\doublespacing
\setlength{\parindent}{0.5in}
\setlength{\parskip}{0pt}
\setlength{\headheight}{15pt}
\pagestyle{fancy}
\fancyhf{}
\fancyhead[R]{\thepage}
\renewcommand{\headrulewidth}{0pt}
\renewcommand{\footrulewidth}{0pt}
\fancypagestyle{plain}{%
  \fancyhf{}
  \fancyhead[R]{\thepage}
  \renewcommand{\headrulewidth}{0pt}
}

\titleformat{\section}
  {\normalfont\bfseries\centering}
  {}{0pt}{}
\titlespacing*{\section}{0pt}{1em}{0.5em}

\newcommand{\referenceentry}[1]{%
  \par\noindent\hangindent=0.5in\hangafter=1 #1\par
}

\begin{document}

\begin{titlepage}
\thispagestyle{plain}
\begin{center}
\vspace*{1.5in}
\textbf{Critical Reviews on Health Across the Lifespan}\\
\vspace{1.5em}
Zixi Li\\
Faculty of Health, Southern Cross University\\
NURS5009-2026-T4 The Health Narrative Across the Lifespan\\
Mrs Debbie Proctor\\
October 3, 2026
\end{center}
\end{titlepage}

\setcounter{page}{2}
\begin{center}
\textbf{Critical Reflections on Health Across the Lifespan}
\end{center}

\section*{Module 3: Adolescent Mental Health and Suicide Risk}

Calear et al. (2026) examined multilevel risk and protective factors for self-harm, suicidal ideation and suicide attempts in Australian adolescents. The prospective study with 4,122 students in 134 schools found that mental health difficulties, peer problems, female or gender-diverse identities and LGBTQA+ identities were associated with self-harm or suicidal ideation. Peer problems were also associated with suicide attempts. Unexpectedly, greater prosocial behaviour was positively associated with self-harm and suicidal ideation. The large Australian cohort and longitudinal methodology was a strength of the study, which investigated risks over a 12-month time frame across individual, family, social and school influences.

However, these results need to be viewed with caution. Attrition was higher among those adolescents with worse mental health and suicidality; only 4,122 of the 5,720 students who completed the baseline assessment participated at follow-up, and this could have resulted in an underestimation of risk. Also, measures used were primarily adolescent self-report, and this was an observational study, so it showed correlation rather than causality. The counterintuitive finding that greater prosocial behaviour was associated with self-harm and suicidal ideation was not significant on univariate analysis, and so its association may reflect complex relationships with other variables in the multivariate model. The study can be used to guide early identification priorities, but cannot demonstrate that implementation of a particular screening or nursing intervention would be effective in reducing suicide attempts. It also provides limited insight into how nurses should respond post-identification.

This article connects with ULO 1 because it shows how health and wellbeing are shaped by psychological, social and environmental factors. It also relates to ULO 3 by promoting the awareness of respectful, inclusive assessment for gender-diverse and LGBTQA+ adolescents, free of prejudice or stigma. ULO 4 is evident as it considers adolescence as a phase of development, with changing emotional, social and cognitive needs. 

For nursing practice, this article underscores the importance of non-judgemental communication and holistic psychosocial assessment, including peer interactions, bullying, family support and protection factors.  I would employ structured assessment to facilitate early recognition, but not rely on a checklist alone. Safe planning, referral and follow-up are vital.

\section*{Module 4: Neighbourhood Context and Psychological Distress}

Learnihan et al. (2026) found the more socially fragmented the neighbourhood, the higher the mean psychological distress score was among Brisbane adults aged 40--65 years. Based on longitudinal HABITAT data from 2,902 men and 3,950 women linked to the census, and controlling for covariates, male residents in the most socially fragmented quintile had a 0.66 point higher score (out of 24) than those in the least quintile, whereas the result was not statistically significant for women after controlling for neighbourhood disadvantage.

Precaution needed, since the small difference could be significant for population planning but might be of limited value for individual clinical diagnosis and prediction. The area index measures something related to a neighbourhood but not necessarily individual loneliness or belonging. This is observational data without the possibility of establishing cause and effect; the analysis was limited to participants who remained residentially stable for at least seven years, which may reduce generalisability to more mobile populations; and the Kessler-6 measures mainly internalising symptoms and may underestimate distressed presentation in men.

The article considers ULO 1: social environments influence health in conjunction with personal factors; ULO 3: nurses should recognise that distress experienced in a community context reflects structural and cultural conditions, not personal failing, and should approach assessment with cultural responsiveness and respect for diverse identities; and ULO 4: at midlife, balancing the responsibilities of caregiving and work life with local resources. In practice, this article and research encourage nurses to screen for social connection, community capital, and physical barriers in health checks.

I believe that these findings suggest that the living environment may contribute to wellbeing alongside individual factors. I realise that while neighbourhood-level data cannot replace individual assessment, awareness of community context can support more holistic, person-centred care and inform advocacy for community resources.

\section*{Module 5: Social Media, Hearing Loss and Social Connection}

Jayakody et al. (2022) identified the perspectives of Australian older people regarding the development of a Facebook community to promote psychosocial wellbeing of hearing-impaired older people. They held a community conversation with 40 people aged 60 years or older to gain input into the design and contents of an online community. Most participants did not use Facebook and were not interested in joining such a group, and favoured phone calls and in-person activities. They identified concerns about lack of access to, and confidence with, the internet as well as privacy and usability.

However, the results should not be taken as evidence for or against social media's effect on loneliness. The study examined opinions about a hypothetical Facebook group but did not create the group or assess loneliness, social isolation or wellbeing outcomes. Only 4 of the 40 participants occasionally used Facebook, and the sample may not be representative of all Australian older people or those with hearing loss, particularly those who already use digital platforms and may be less socially isolated. Group discussions may have led people to "go along" with others, and the article has limited information on whether participants' views were influenced by age, digital literacy or hearing needs. Therefore, the results identify barriers and preferences but do not establish whether alternative strategies are more effective.

This article relates to ULO 1 as it separates objective social contact from perceptions of contact and wellbeing. It relates to ULO 3 through respect for the older person's preferences, independence and dignity; participants noted that certain forms of address could be perceived as condescending. It relates to ULO 4 through the characterisation of the influences of hearing loss and ageing on access to communication and social participation. 

For use in nursing the article serves to discourage all internet or digital connection prescriptions. I would enquire into the older person's most accessible and desirable contact forms and take into consideration logistical challenges such as affordable tech support and internet, hearing aids and privacy, safety, competence and confidence to use digital technology. Phone calls and local, face-to-face community groups may be more acceptable options to some people. Digital options should be offered alongside, not instead of, approaches chosen with the older person.

\newpage
\section*{References}

\referenceentry{Calear, A. L., Batterham, P. J., Werner-Seidler, A., Maston, K., Torok, M., O'Dea, B., Larsen, M. E., \& Christensen, H. (2026). Multilevel risk and protective factors for self-harm, suicidal ideation and suicide attempt in adolescents. \textit{Journal of Child Psychology and Psychiatry, 67}(5), 609--619. \url{https://doi.org/10.1111/jcpp.70024}}

\referenceentry{Jayakody, D. M. P., Tan, Y. M. E., Livings, I., Costello, L., Flicker, L., \& Almeida, O. P. (2022). Australian older adults' views on using social media for reducing social isolation and loneliness in hearing impaired older adults: A community conversation. \textit{Australasian Journal on Ageing, 41}(4), 585--589. \url{https://doi.org/10.1111/ajag.13137}}

\referenceentry{Learnihan, V., Bagheri, N., \& Turrell, G. (2026). The longitudinal association between living in socially fragmented neighbourhoods and psychological distress among middle aged adults in Brisbane, Australia. \textit{Journal of Affective Disorders, 401}, Article 121318. \url{https://doi.org/10.1016/j.jad.2026.121318}}

\newpage
\section*{GenAI Acknowledgement Statement}
\noindent I did not use generative AI in the preparation of this assessment. I confirm that the content of this critical reflection is entirely my own work, and no part of it was generated, rewritten, or substantially improved by generative artificial intelligence tools.

\end{document}
